\begin{flushright}
  \fechaclase{Clase del 25 de agosto de 2026}
\end{flushright}

\paragraph{Control no lineal}
El control lineal es un campo bien establecido con numerosos métodos efectivos
y una larga trayectoria de aplicaciones exitosas. Sin embargo, el control no
lineal ha crecido en interés en áreas como robótica, biomédica, control de
aeronaves, etc. Pues

\begin{itemize}
  \item Mejoran los sistemas de control existentes en comparación con el
    control lineal
  \item Manejo de incertidumbres en el modelo.
  \item Simplicidad de diseño: los controladores no lineales pueden ser más
    fáciles de diseñar, al estar basados en la física del sistema
\end{itemize}

El control no lineal es una disciplina clave en la ingeniería de control, que
permite abordar problemas prácticos más complejos y ofrece un entendimiento
más profundo de los sistemas reales.

\paragraph{Control robusto}
Es una disciplina dentro de la teoría de control que se enfoca en el diseño de
sistemas de control capaces de mantener su rendimiento deseado a pesar de
incertidumbres y perturbaciones.

Las técnicas más comunes son el ADRC, SMC y algunas variantes de
\textit{backstepping}.

\subsection{Preámbulo matemático}

\paragraph{Kernel o espacio nulo de una transformación lineal}
\[
  A\colon \mathbb{C}^{n}\longrightarrow\mathbb{C}^{m}
\]
se define como
\[
  \operatorname{Ker}(A)=N(A)
  :=\left\{x\in\mathbb{C}^{m}\mid Ax=0\right\}
\]

\paragraph{La imagen o rango de una transformación lineal se define como}
\[
  \operatorname{Im}(A)=R(A)
  :=\left\{y:\mathbb{C}^{m}\mid Ax=y,\quad
  x\in\mathbb{C}^{n}\right\}
\]

\paragraph{Potencias de matrices:} para un número no negativo entero $p$
\[
  A^{p}=\underbrace{AA\cdots A}_{p\text{ veces}},
  \qquad A^{0}=I
\]

\paragraph{Las leyes de los exponentes se cumplen}
\[
  A^{p}A^{q}=A^{p+q},
  \qquad
  (A^{p})^{q}=A^{pq}
\]

\paragraph{Si las matrices cumplen $AB=BA$ entonces}
\[
  (AB)^{p}=A^{p}B^{p}
\]

\paragraph{Transpuesta de dos matrices:}
\[
  (AB)^{T}=B^{T}A^{T}
\]
pues
\[
  (AB)^{T}
  =\left[\sum_{k=1}^{h}a_{jk}b_{ki}\right]_{i,j=1}^{m,p}
  =\left[\sum_{k=1}^{h}b_{ki}a_{jk}\right]_{i,j=1}^{m,p}
  =B^{T}A^{T}
\]

\paragraph{Matriz simétrica:}
Se dice que una matriz cuadrada $A\in\mathbb{R}^{n\times n}$ es simétrica si
cumple
\[
  A=A^{T}
\]

\paragraph{Matriz no singular:}
Se dice que la matriz cuadrada $A\in\mathbb{R}^{n\times n}$ es no singular si
\[
  \det(A)\neq 0
\]
en otro caso se dice que es singular.

Además
\[
  (A^{-1})^{T}=(A^{T})^{-1}
\]
pues
\[
  I_{n\times n}=(AA^{-1})^{T}=(A^{-1})^{T}A^{T}
\]